\documentclass[11pt,a4paper]{article}
\usepackage[utf8]{inputenc}
\usepackage{amsmath,amsfonts,amssymb}
\usepackage{geometry}
\geometry{margin=1in}
\usepackage{hyperref}
\usepackage[many]{tcolorbox}
\usepackage{enumitem}
\usepackage{fancyhdr}

\pagestyle{fancy}
\fancyhf{}
\fancyfoot[L]{Universitas Bengkulu - Teknik Elektro}
\fancyfoot[R]{\thepage}
\def\footrule{{\color{gray}\hrule}}

\title{\vspace{-2cm}\textbf{Problem Set (Tugas Terstruktur)\\Persamaan Diferensial - Minggu 1: Klasifikasi \& Pemodelan}}
\author{Novalio Daratha \& Muhammad Arfan\\Program Studi Teknik Elektro\\Universitas Bengkulu}
\date{2026}

\begin{document}

\maketitle

\begin{tcolorbox}[colback=red!5!white,colframe=red!75!black,title=Instruksi Tugas]
Kumpulkan jawaban Problem Set ini pada awal perkuliahan minggu berikutnya. Tunjukkan semua langkah pemikiran, analisis fisis, dan penurunan matematis Anda. Serta lampirkan bukti validasi jika diinstruksikan.
\end{tcolorbox}

\section*{Daftar Soal (Level Kognitif C1 - C6)}

\begin{enumerate}[label=\textbf{\arabic*.}]
    \item \textbf{(C1 - Mengingat)} Tuliskan definisi lengkap dari Persamaan Diferensial Biasa (PDB) linier orde 2. Berikan satu contoh persamaan yang memenuhi kriteria tersebut di bidang kelistrikan.
    \item \textbf{(C2 - Memahami)} Jelaskan dengan bahasa Anda sendiri perbedaan mendasar antara "Solusi Umum" dan "Solusi Khusus" pada penyelesaian suatu \textit{Initial Value Problem} (IVP). Mengapa kita membutuhkan kondisi awal?
    \item \textbf{(C3 - Mengaplikasikan)} Tentukan apakah persamaan-persamaan diferensial berikut merupakan persamaan linier atau non-linier, serta sebutkan ordenya:
    \begin{enumerate}[label=\alph*.]
        \item $\frac{d^2y}{dt^2} + 3t\frac{dy}{dt} + y = \sin(t)$
        \item $\left(\frac{dx}{dt}\right)^2 + 4x = 0$
        \item $\frac{\partial^2 V}{\partial x^2} + \frac{\partial^2 V}{\partial y^2} = \rho(x,y)$
    \end{enumerate}
    \item \textbf{(C4 - Menganalisis)} Analisislah sebuah rangkaian seri RC (Resistor-Kapasitor) yang dihubungkan dengan sumber tegangan DC sebesar $V_0$. Dengan menggunakan Hukum Tegangan Kirchhoff (KVL), turunkan pemodelan fisis rangkaian tersebut hingga membentuk PDB orde 1 terhadap variabel muatan kapasitor $q(t)$.
    \item \textbf{(C5 - Mengevaluasi)} Diberikan model PDB rangkaian listrik: $\frac{di}{dt} + 5i = 10$ dengan kondisi awal $i(0)=0$. Seorang praktikan mengklaim solusi khususnya adalah $i(t) = 2 - 2e^{-5t}$. Evaluasi dan buktikan apakah klaim tersebut benar atau salah!
    \item \textbf{(C6 - Mencipta)} Rancanglah sebuah sistem elektromekanik atau sirkuit listrik imajiner sederhana sedemikian rupa sehingga dinamika responsnya direpresentasikan oleh $\frac{dx}{dt} + kx = F_0$. Jelaskan secara detil komponen-komponen sistem Anda, variabel $x$, serta makna fisis dari parameter $k$ dan $F_0$. \textbf{Tugas Tambahan:} Tulis fungsi Python sederhana (Euler Method) di secarik kertas untuk menyelesaikan persamaan yang Anda rancang.
\end{enumerate}

\end{document}
